\chapter{Kết luận}\label{chapter-7-conclusion}

Khóa luận này tìm cách tăng tốc test-time discovery trong sinh mã phức tạp bằng execution-guided self-distillation. Đóng góp phương pháp luận chính là một tổ chức teacher-first: self-teacher sinh các quỹ đạo dưới sự dẫn dắt của execution feedback, rồi một verifier và một semantic judge lọc chúng theo tính đúng chức năng và tính độc lập cấu trúc. Bộ lọc verifier--judge đảm bảo một lời giải sao chép không bao giờ lọt vào diện được distill. Cách tổ chức này đặt phương pháp đề xuất ở đâu đó giữa Self-Distillation Policy Optimization (SDPO) on-policy và Supervised Fine-Tuning (SFT) off-policy.

Đánh giá thực nghiệm cho ra một hồ sơ hiệu năng compute-fair. Trên một đánh giá matched quy mô trung bình, teacher-first vượt trội rõ rệt so với baseline student-first, và thoát được flat-reward trap trên chính những bài lập trình khó nhất. Dưới một ngân sách generation cân bằng, và qua phân tích compute-to-correct frontier, lợi thế này vẫn giữ vững một cách nhất quán — cho thấy mức tăng là một thuộc tính thuật toán của quỹ đạo distillation được chọn, chứ không phải artifact của khối lượng sampling. Nghiên cứu cũng nhận thấy cách formulate execution feedback qua reprompt-template có ảnh hưởng đến discovery, rõ nhất trên các bài khó.

Khi bật reasoning trace, distill từ một feedback-conditioned context làm giảm epistemic verbalization một cách đo được trong sinh mã, và ý nghĩa của sự đè nén này lại phụ thuộc vào regime. Điều đáng chú ý nhất: với cùng một mô hình dùng trên cả hai miền, tương phản code-versus-math vạch ra một ranh giới value-versus-procedure khá rõ. Trên code — nơi output là một phương pháp mà teacher có thể cài đặt được — mô hình thoát ra khỏi trap; còn trên toán chỉ-có-đáp-số — nơi reference chỉ là một giá trị trơ mà teacher chỉ có thể sao chép — mô hình kẹt nguyên ở mức 0 trên pilot AIME. Điều này khá nhất quán với cách đọc rằng một sự thoát thành công tại thời điểm suy luận đòi hỏi privileged reference phải mã hóa một phương pháp trong tầm với, chứ không chỉ một giá trị cuối cùng trơ trọi — dù phía toán vẫn chỉ dựa trên một pilot nhỏ (n = 2) nên được đọc theo hướng directional.

Công trình này vẫn là một đặc trưng hóa thực nghiệm trên một mô hình 4B duy nhất, trong phạm vi compute cá nhân, và việc mở rộng lên các kiến trúc mạnh hơn nhiều khả năng sẽ thu hẹp biên so với baseline. Trong phạm vi đó, nó đưa ra một mô tả compute-fair, giải thích theo hướng cơ chế về execution-guided self-distillation, cùng một đặc trưng hóa cho ranh giới value-versus-procedure. Bằng cách vạch ra khi nào và bằng cách nào cách tiếp cận này tăng tốc được test-time discovery, nó đặt một nền tảng cho các mở rộng tương lai trong tối ưu tại thời điểm suy luận.
